\begin{lemma}
Let $G$ be a finite graph, let $\Delta$ be a positive integer, and let
$\gamma>0$.  If the activation--priority sampling law
\noderef{RandomIndependentSetSampling} exists on a nonempty vertex set with
activation probability $\gamma/\Delta$, then $\gamma\le \Delta$.
\end{lemma}
\begin{proof}
Let $A=V(G)$ be the atom in which every vertex is activated.  By
\noderef{RandomIndependentSetSampling}, the measure of this atom is
\[
\left(\frac{\gamma}{\Delta}\right)^{|V(G)|}
\left(1-\frac{\gamma}{\Delta}\right)^0
=\left(\frac{\gamma}{\Delta}\right)^{|V(G)|}.
\]
The atom is a subset of the whole sample space, whose total measure is $1$,
so this number is at most $1$.  Since the vertex set is nonempty, its
cardinality is positive.  If $\gamma/\Delta>1$, then the positive power
$(\gamma/\Delta)^{|V(G)|}$ would be strictly larger than $1$, contradicting
the preceding inequality.  Therefore $\gamma/\Delta\le 1$.  Because
$\Delta>0$, multiplying by $\Delta$ gives $\gamma\le\Delta$.
\end{proof}
